\chapter{Introduction}

The field of high energy physics is the study of all fundamental particles and their interactions at extremely short distances, which are described by a quantum field theory known as the Standard Model of particle physics. The Standard Model is an extremely successful framework which predicts and calculates quantities that can be measured in high energy physics experiments. Perhaps one of the most notable predictions of the Standard Model is that of the top quark. This particle was predicted long before it was discovered in 1995 because the Standard Model required a particle with its charge and intrinsic spin to conserve several of the fundamental symmetries built into the model.

The top quark is unique because of its large mass, which is nearly that of a gold atom. Because of its large mass, creating top quarks requires a tremendously energetic particle accelerator. Such an accelerator exists, called the Tevatron, located at the Fermi National Laboratory, which collides protons ($p$) and antiprotons ($\bar{p}$) at a center of mass energy of 1.96~TeV. An aerial view of the accelerator facility is shown in Figure~\ref{Fermilab}. At the Tevatron the top quark is primarily studied when a top ($t$) and antitop ($\bar{t}$) quark are produced together;  this process occurs for roughly one out of every five billion $\ppbar$ collisions. The Standard Model also predicts that the top quark can be produced singly at the Tevatron; however, it is expected to occur at half the rate of top-antitop production, or in nearly one out of every ten billion collisions. 

Measuring single top quark production is interesting for several reasons, but perhaps the most compelling reason is that one can test the unitarity of the matrix ($V_{CKM}$) that governs how quarks mix. If unitarity is violated (i.e. $VV^{\dagger}\neq I$) it would provide clear evidence for physics beyond the Standard Model such as a possible fourth quark generation. Models which predict new processes that are not allowed by the Standard Model can also be tested by measuring the single top production rate. For instance, many theories beyond the Standard Model predict a large ``flavor changing neutral current'' rate which could enhance the single top production rate by one order of magnitude. Finally, single top is important to measure because it represents one of the largest irreducible backgrounds for Higgs production at the Tevatron. The Higgs particle is predicted by the theory that explains why fundamental particles have mass and remains the only Standard Model particle not to have been found by high energy experiments.

\begin{figure}[!h!tbp]
\begin{center}
\includegraphics[width=0.75\textwidth]{eps/Tevatron/TevatronArial.eps}
\end{center}
\vspace{-0.1in}
\caption{Aerial view of the Fermilab accelerator complex. }
\label{Fermilab}
\end{figure}

An outline of this thesis is as follows. Chapter~\ref{theory} of this thesis gives an introduction to the Standard Model of particle physics. Also described in this chapter is the motivation for measuring properties of the top quark along with a detailed description of single top quark production. Chapter~\ref{experiment} describes the Tevatron particle accelerator at Fermilab used to create $\ppbar$~collisions at $\sqrt{s}$=1.96~TeV. The decay products of the proton-antiproton collisions are measured using the $\dzero$~detector, which is also described in this chapter. Chapter~\ref{recosim} explains how signals in the $\dzero$~detector are used to reconstruct the collision to determine if it is consistent with single top quark production.

Chapter~\ref{data} describes the Tevatron Run~II dataset used to search for single top quark production. Because single top quark production is rare and its signature is common to many other physics processes a proper modeling of the background is crucial to the analysis. Chapter~\ref{backgrounds} describes the main sources of background to single top quark production as well as how they are modeled.

To better identify the single top signal contribution in the dataset a technique known as the matrix element method is employed. This technique uses leading order matrix elements to create an event probability for both signal and background hypotheses. These probabilities are combined to form a variable for each event which is designed to peak near one for signal-like events or near zero for background-like events. This technique and its application is described in chapter~\ref{matrixelement}.

Chapter~\ref{limits} describes the sophisticated Bayesian statistical analysis used to extract the single top quark cross section in the dataset. A detailed discussion of the systematic errors in the signal and background modeling is also provided here. Chapter~\ref{results} shows the results of the single top quark analysis including the measured cross section and signal significance. Finally, chapter~\ref{conclusion} concludes the thesis with a few remarks on future improvements to the analysis technique as well as possible measurements with an increased dataset.